\documentclass[11pt,a4paper]{article}

\usepackage[utf8]{inputenc}
\usepackage[T1]{fontenc}
\usepackage{lmodern}
\usepackage[a4paper,margin=1.9cm]{geometry}
\usepackage{xcolor}
\usepackage{enumitem}
\usepackage{graphicx}
\usepackage{parskip}

\definecolor{accent}{HTML}{1F3A93}
\definecolor{band}{HTML}{EAEEF6}
\pagestyle{empty}
\setlength{\parskip}{5pt}

\newcommand{\heading}[1]{\par\medskip{\color{accent}\bfseries\large #1}\par\nopagebreak}

\begin{document}

% ----------------------------------------------------------------- header
\begin{center}
{\color{accent}\Large\bfseries Decarbonise Your Equity Portfolio --
Without Giving Up Performance}\\[3pt]
{\large Group BR Asset Management \quad$\vert$\quad Net-Zero Equity Strategy}
\end{center}

\noindent\textcolor{accent}{\rule{\textwidth}{1.2pt}}

% ----------------------------------------------------------------- objective
\heading{The objective}

As a pension fund, you have made a net-zero commitment to your beneficiaries and
your regulator. You also have a fiduciary duty to deliver competitive long-term
returns. Those two goals are usually presented as a trade-off. They are not. Our
\textbf{Net-Zero Equity Strategy} delivers a credible, year-on-year
decarbonisation path on your North American and European equity allocation while
keeping the portfolio's risk and return statistically indistinguishable from
your current benchmark.

% ----------------------------------------------------------------- strategy
\heading{The strategy in plain terms}

We start from your existing value-weighted benchmark and keep it almost intact.
Each year we make one targeted adjustment: we trim the small group of
carbon-heavy companies -- mainly regulated utilities and a few cement and oil
producers -- that generate the bulk of the portfolio's emissions, and we
redeploy that capital across the hundreds of cleaner companies you already own.
The portfolio stays fully invested, long-only, and broadly diversified (around
1{,}000 holdings). A transparent optimiser guarantees that the carbon budget is
met every year while the portfolio stays as close as mathematically possible to
your benchmark. The carbon ceiling tightens by \textbf{10\% per year}, putting
you on an explicit net-zero glide path.

% ----------------------------------------------------------------- evidence
\heading{The evidence (back-test 2014--2024, 11 years)}

\begin{center}
\colorbox{band}{\begin{minipage}{0.95\textwidth}\centering\vspace{4pt}
\begin{tabular}{p{0.30\textwidth} p{0.62\textwidth}}
\textbf{$-$50\% carbon footprint} &
on a one-off decarbonised mandate -- and a continuing 10\%-a-year reduction
under the net-zero glide path.\\[4pt]
\textbf{$-$0.03 Sharpe ratio} &
the only performance cost: 0.69 versus 0.72 for the benchmark. Annual return
moved from 10.6\% to 10.2\%.\\[4pt]
\textbf{1.1\% tracking error} &
the decarbonised portfolio shadows your benchmark almost perfectly -- no style
drift, no hidden bets.
\end{tabular}
\vspace{4pt}\end{minipage}}
\end{center}

In short: we removed half of the portfolio's emissions for a performance
difference smaller than the noise in a single quarter. We also found that
``defensive'' low-volatility strategies are, in this market, a hidden
\emph{carbon} bet -- a minimum-variance portfolio carried five times the
benchmark's footprint -- so screening for carbon is also a useful risk control.

% ----------------------------------------------------------------- risks
\heading{Trade-offs and risks -- and how we manage them}

\begin{itemize}[nosep,leftmargin=1.4em]
  \item \textbf{Sector tilt.} The strategy is structurally under-weight
        utilities and energy. We cap that tilt and minimise tracking error so
        the portfolio cannot drift far from the benchmark in any single year.
  \item \textbf{Regime risk.} Decarbonisation was nearly free over the last
        decade partly because clean leaders outperformed. In a fossil-led
        market the carbon screen could cost return; the 1.1\% tracking-error
        budget bounds that downside, and the glide path can be paused.
  \item \textbf{Feasibility.} A long-only portfolio cannot decarbonise
        indefinitely. We monitor the remaining carbon headroom each year and
        will report transparently when the glide path approaches its
        mathematical limit, well before it binds.
  \item \textbf{Data quality.} Emissions data are self-reported and lag by a
        year. We use Scope 1+2 figures, forward-fill conservatively, and update
        the carbon budget annually as disclosure improves.
\end{itemize}

\vspace{2pt}
\noindent\textcolor{accent}{\rule{\textwidth}{1.2pt}}
\begin{center}
\textbf{Halve your equity carbon footprint today, reach net zero on a clear
glide path -- and keep the portfolio you already trust.}
\end{center}

\end{document}
